\documentclass[12pt,a4paper]{article}

\usepackage{amsmath, amssymb, graphicx}
\usepackage{geometry}
\usepackage{setspace}
\usepackage{hyperref}

\geometry{margin=1in}
\setstretch{1.2}

\title{\textbf{Lecture Scribe}\\Fundamentals of Probability in Computing}
\author{
Name: Sakshi Patel\\
Student ID: AU2440206\\
Course: Fundamentals of Probability in Computing\\
}
\date{}
\begin{document}

\maketitle
\newpage

\section*{1️⃣ Main Topics Covered in the Lecture}

From the lecture material, the key concepts are:

\begin{itemize}
\item Gaussian Random Variable
\item Q-function and Error Function
\item Relation between $\Phi$ and $Q$
\item Continuous Random Variables
\begin{itemize}
\item Uniform
\item Exponential
\item Laplace
\item Gamma
\end{itemize}
\item Graphs \& special cases
\item Density estimation \& problem solving
\end{itemize}

We’ll now build understanding in a logical teaching order.

\section*{2️⃣ Continuous Random Variables — Foundation}

\textbf{Conceptual idea}

A variable that can take any real value within an interval.

Examples:

\begin{itemize}
\item Time taken by an algorithm
\item Network latency
\item Temperature
\end{itemize}

\textbf{Mathematical idea}

Probability at a single point = 0

\[
P(X = x) = 0
\]

Instead:

\[
P(a < X < b) = \int_a^b f(x)\,dx
\]

where $f(x)$ = Probability Density Function (PDF)

\textbf{Intuition}

Probability = area under curve.

Why?

Because infinite possible values → probability spread continuously.

\section*{3️⃣ Gaussian Random Variable (Normal Distribution)}

\textbf{Definition}

A continuous random variable whose values cluster around a mean following a bell-shaped curve.

\textbf{Formula}

\[
f(x) = \frac{1}{\sqrt{2\pi\sigma^2}} e^{-\frac{(x-\mu)^2}{2\sigma^2}}
\]

Where:

\begin{itemize}
\item $\mu$ → mean
\item $\sigma^2$ → variance
\end{itemize}

\textbf{Conceptual Understanding}

Nature loves Gaussian distributions:

\begin{itemize}
\item Human height
\item Sensor noise
\item Measurement errors
\end{itemize}

Reason:

Many small random effects combine → normal distribution emerges.

\textbf{Why this formula? (Stepwise reasoning)}

We need:

\begin{itemize}
\item Symmetry around mean → depends on $x-\mu$
\item Equal effect left/right → square term $(x-\mu)^2$
\item Smooth decay → exponential function
\item Total probability = 1 → normalization constant $\frac{1}{\sqrt{2\pi\sigma^2}}$
\end{itemize}

Each step logically builds:

symmetry → squared distance → exponential decay → normalization.

\textbf{Properties}

\begin{itemize}
\item Symmetric curve
\item Mean = median = mode
\item Bell shaped
\item Controlled spread by variance
\end{itemize}

\textbf{Example}

Marks distribution:

Mean = 60

Std dev = 10

Probability student scores between 50 and 70:

Step 1: Convert to standard normal

\[
Z = \frac{X - \mu}{\sigma}
\]

\[
Z_1 = -1, \quad Z_2 = 1
\]

Step 2: Use Z-table

\[
P(-1 < Z < 1)
\]

Step 3: Interpret area.

\section*{4️⃣ Standard Normal Distribution}

When:

\[
Z = \frac{X - \mu}{\sigma}
\]

Then:

\begin{itemize}
\item Mean = 0
\item Variance = 1
\end{itemize}

Why needed?

Because every Gaussian becomes one standard form → easy calculations.

\section*{5️⃣ $\Phi$ (Phi) Function}

\[
\Phi(z) = P(Z \le z)
\]

Represents cumulative probability from left side.

Conceptually:

Area under curve till point $z$.

\section*{6️⃣ Q-Function}

\[
Q(z) = P(Z > z)
\]

Right-tail probability.

Used in:

\begin{itemize}
\item Communication systems
\item Error probability
\item Machine learning thresholds
\end{itemize}

\textbf{Relation between $\Phi$ and $Q$}

\[
Q(z) = 1 - \Phi(z)
\]

Logical reasoning:

Total probability = 1

Left area + right area = 1.

\section*{7️⃣ Error Function}

Used for Gaussian integration.

\[
\text{erf}(x) = \frac{2}{\sqrt{\pi}} \int_0^x e^{-t^2} dt
\]

Relation:

\[
\Phi(z) = \frac{1}{2}\left[1 + \text{erf}\left(\frac{z}{\sqrt{2}}\right)\right]
\]

Why used?

Because Gaussian integrals don’t have simple closed forms.

\section*{8️⃣ Types of Continuous Random Variables}

\subsection*{�� Uniform Distribution}

Definition

All values equally likely.

\[
f(x) = \frac{1}{b-a}, \quad a \le x \le b
\]

Intuition

Flat probability.

Example

Random number generator between 0 and 1.

\subsection*{�� Exponential Distribution}

Definition

Models waiting time for events.

\[
f(x) = \lambda e^{-\lambda x}
\]

Property — Memoryless

\[
P(X > s+t \mid X > s) = P(X > t)
\]

Why?

Future independent of past.

Example

Time until server failure.

\subsection*{�� Laplace Distribution}

Double exponential.

\[
f(x) = \frac{1}{2b} e^{-|x-\mu|/b}
\]

Sharper peak than Gaussian.

Used in:

\begin{itemize}
\item Signal compression
\item Data modelling with outliers
\end{itemize}

\subsection*{�� Gamma Distribution}

Generalization of exponential.

\[
f(x) = \frac{x^{k-1} e^{-x/\theta}}{\Gamma(k)\theta^k}
\]

Special case

$k = 1 \rightarrow$ exponential.

Example

Time until $k$ events occur.

\section*{9️⃣ Graphical Understanding}

\begin{itemize}
\item Uniform → flat
\item Gaussian → bell
\item Exponential → decreasing
\item Laplace → sharp center
\item Gamma → flexible skewed shape
\end{itemize}

Each distribution models different real-world uncertainty.

\section*{10 Gaussian Density Estimation}

Goal:

Find distribution parameters from data.

Steps:

\begin{itemize}
\item Collect data
\item Calculate mean
\item Calculate variance
\item Fit Gaussian curve
\end{itemize}

Used in:

\begin{itemize}
\item Machine learning
\item Pattern recognition
\item Signal processing
\end{itemize}

\section*{1️⃣1️⃣ Problem Solving Logic}

Whenever a probability question appears:

\begin{itemize}
\item Step 1 → Identify distribution
\item Step 2 → Write PDF
\item Step 3 → Convert to standard form if Gaussian
\item Step 4 → Integrate / use table
\item Step 5 → Interpret result
\end{itemize}

\section*{1️⃣2️⃣ Concept Check}

\begin{tabular}{|l|l|l|l|}
\hline
Concept & Mathematical & Intuitive & Where used\\
\hline
Gaussian & exponential squared & natural randomness & ML, signals\\
Uniform & constant & equal chance & simulations\\
Exponential & decay & waiting time & reliability\\
Gamma & generalized & multiple events & queueing\\
\hline
\end{tabular}
\newline\\
Yes — these interpretations match real problems.

\section*{13 Concise Summary}

\begin{itemize}
\item Continuous variables → probability via area
\item Gaussian → most important real-world distribution
\item $\Phi$ → cumulative probability
\item $Q$ → error/tail probability
\item Uniform, exponential, Laplace, gamma → key modeling tools
\item Density estimation → find distribution from data
\end{itemize}

\section*{14 Key Formula Sheet}

Gaussian

\[
f(x) = \frac{1}{\sqrt{2\pi\sigma^2}} e^{-(x-\mu)^2/2\sigma^2}
\]

Standardization

\[
Z = \frac{X-\mu}{\sigma}
\]

$\Phi$-function

\[
\Phi(z) = P(Z \le z)
\]

Q-function

\[
Q(z) = 1-\Phi(z)
\]

Uniform

\[
f(x)=\frac{1}{b-a}
\]

Exponential

\[
f(x)=\lambda e^{-\lambda x}
\]

Laplace

\[
f(x)=\frac{1}{2b} e^{-|x-\mu|/b}
\]

Gamma

\[
f(x)=\frac{x^{k-1} e^{-x/\theta}}{\Gamma(k)\theta^k}
\]

\section*{15 Probable Exam Questions}

\textbf{Theory}

\begin{itemize}
\item Define Gaussian distribution
\item Properties of exponential distribution
\item Laplace vs Gaussian
\end{itemize}

\textbf{Derivations}

\begin{itemize}
\item Standard normal transformation
\item $\Phi$-$Q$ relation
\item Mean/variance of exponential
\end{itemize}

\textbf{Numerical}

\begin{itemize}
\item Z-table probability
\item Uniform distribution interval
\item Waiting time problems
\end{itemize}

\textbf{Conceptual}

\begin{itemize}
\item Why Gaussian appears naturally
\item Memoryless property explanation
\end{itemize}

\end{document}
